\documentclass[10pt,xcolor={dvipsnames}]{beamer}

% Tema profesional estándar
\usetheme{Copenhagen}
\usecolortheme{seahorse}

% Personalización de colores
\setbeamercolor{structure}{fg=NavyBlue}
\setbeamercolor{palette primary}{bg=NavyBlue,fg=white}
\setbeamercolor{palette secondary}{bg=NavyBlue!80,fg=white}
\setbeamercolor{palette tertiary}{bg=NavyBlue!60,fg=white}
\setbeamercolor{frametitle}{bg=NavyBlue,fg=white}
\setbeamercolor{block title}{bg=NavyBlue!80,fg=white}
\setbeamercolor{block body}{bg=NavyBlue!10,fg=black}

% Configuración adicional
\setbeamertemplate{navigation symbols}{}
\setbeamertemplate{footline}[frame number] 

%-------------------------------------------------------
% INCLUDE PACKAGES
%-------------------------------------------------------

\usepackage[utf8]{inputenc}
\usepackage[spanish,es-tabla]{babel}
\usepackage[T1]{fontenc}
\usepackage{graphicx}
\graphicspath{{img/}}
\usepackage{amsmath}
\usepackage{amssymb}

%-------------------------------------------------------
% DEFFINING AND REDEFINING COMMANDS
%-------------------------------------------------------

% colored hyperlinks
\usepackage{hyperref}
\hypersetup{
    colorlinks=true,
    linkcolor=NavyBlue,
    urlcolor=NavyBlue
}

%-------------------------------------------------------
% INFORMATION IN THE TITLE PAGE
%-------------------------------------------------------

\title[] % [] is optional - is placed on the bottom of the sidebar on every slide
{ % is placed on the title page
      \textbf{Diseño e Implementación de una Plataforma de Estabilización para una Bicicleta Utilizando un Volante de Inercia}
}

\subtitle[Estabilización con Volante de Inercia]
{
      \textbf{Control Giroscópico CMG}
}

\author[Hernández \& Reyes]
{      Jeferson Hernán Hernández Moreno \\
      Fabián Andrés Reyes Romero\\[0.5em]
      {\small Director: Dr. Alexander Jiménez Triana}
}

\institute[]
{%
      Universidad Distrital Francisco José de Caldas\\
      Facultad Tecnológica\\
      Ingeniería en Control y Automatización
}

\date{Enero 2026}

%-------------------------------------------------------
% THE BODY OF THE PRESENTATION
%-------------------------------------------------------

\begin{document}

% ============================================
% PORTADA
% ============================================
\begin{frame}
\titlepage
\end{frame}

% ============================================
% TABLA DE CONTENIDOS
% ============================================
\begin{frame}{Contenido}
\tableofcontents
\end{frame}

% ============================================
% SECCIÓN 1: INTRODUCCIÓN
% ============================================
\section{Introducción}

\begin{frame}{Introducción}
\begin{block}{Contexto}
\begin{itemize}
    \item Búsqueda de soluciones tecnológicas aplicadas a la movilidad
    \item Control de vehículos de dos ruedas en condiciones de baja velocidad
    \item La estabilidad pasiva desaparece a bajas velocidades
\end{itemize}
\end{block}

\begin{block}{Propósito del Proyecto}
Desarrollar una plataforma que aproveche principios físicos de conservación del momento angular para gestionar el equilibrio de la bicicleta en estado de reposo o bajas velocidades.
\end{block}

\vspace{0.3cm}
\begin{center}
\fbox{\parbox{0.8\textwidth}{\centering \textit{[Espacio para imagen: Bicicleta con sistema CMG]}}}
\end{center}
\end{frame}

% ============================================
% SECCIÓN 2: PLANTEAMIENTO DEL PROBLEMA
% ============================================
\section{Planteamiento del Problema}

\begin{frame}{Planteamiento del Problema}
\begin{block}{Problemática en Bogotá}
\begin{itemize}
    \item Congestión del tráfico y movilidad urbana
    \item La bicicleta como alternativa sostenible
    \item Inestabilidad en condiciones de baja velocidad
\end{itemize}
\end{block}

\begin{columns}
\column{0.5\textwidth}
\begin{alertblock}{Riesgos}
\begin{itemize}
    \item Pérdida de equilibrio
    \item Caídas y accidentes
    \item Factor disuasivo para usuarios
\end{itemize}
\end{alertblock}

\column{0.5\textwidth}
\vspace{0.2cm}
\fbox{\parbox{0.9\textwidth}{\centering \textit{[Espacio para gráfica: Estadísticas de accidentes]}}}
\end{columns}

\vspace{0.3cm}
\begin{block}{Vacío Tecnológico}
Ausencia de soluciones locales que proporcionen auto-equilibrio activo a estructuras de dos ruedas.
\end{block}
\end{frame}

% ============================================
% SECCIÓN 3: JUSTIFICACIÓN
% ============================================
\section{Justificación}

\begin{frame}{Justificación}
\begin{block}{Impacto Social}
\begin{itemize}
    \item Reducción de accidentes en biciusuarios
    \item Promoción de movilidad sostenible
    \item Mejora en la calidad del aire en Bogotá
\end{itemize}
\end{block}

\begin{block}{Impacto Tecnológico}
\begin{itemize}
    \item Desarrollo académico y tecnológico nacional
    \item Contribución al estado del arte en sistemas de control
    \item Aplicación práctica de teorías de control avanzado
\end{itemize}
\end{block}

\vspace{0.3cm}
\begin{center}
\fbox{\parbox{0.8\textwidth}{\centering \textit{[Espacio para imagen: Movilidad sostenible]}}}
\end{center}
\end{frame}

% ============================================
% SECCIÓN 4: OBJETIVOS
% ============================================
\section{Objetivos}

\begin{frame}{Objetivo General}
\begin{block}{Objetivo General}
Diseñar y construir una plataforma mecánica experimental, compuesta por una bicicleta equipada con un volante de inercia, sensores, actuadores e instrumentación asociada, que permita implementar, evaluar y validar estrategias de control orientadas a su estabilización.
\end{block}

\vspace{0.3cm}
\begin{center}
\fbox{\parbox{0.8\textwidth}{\centering \textit{[Espacio para diagrama: Sistema completo]}}}
\end{center}
\end{frame}

\begin{frame}{Objetivos Específicos}
\begin{enumerate}
    \item \textbf{Diseño Mecánico}
    \begin{itemize}
        \item Diseñar y construir la estructura mecánica de la plataforma experimental
        \item Acoplar volante de inercia a la bicicleta
    \end{itemize}
    
    \vspace{0.2cm}
    \item \textbf{Integración Electrónica}
    \begin{itemize}
        \item Integrar sensores (IMU), actuadores y sistema de adquisición de datos
        \item Implementar procesamiento en tiempo real
    \end{itemize}
    
    \vspace{0.2cm}
    \item \textbf{Control}
    \begin{itemize}
        \item Implementar estrategia de control PID en lazo cerrado
        \item Validar respuesta del sistema ante perturbaciones
    \end{itemize}
\end{enumerate}
\end{frame}

% ============================================
% SECCIÓN 5: MARCO TEÓRICO
% ============================================
\section{Marco Teórico}

\begin{frame}{Antecedentes - Desarrollos Nacionales}
\begin{block}{Universidad Nacional de Colombia (Baquero Suárez, 2014)}
\begin{itemize}
    \item Control por dirección (manubrio)
    \item Control ADRC-GPI
    \item Limitación: Requiere velocidad mínima de avance
\end{itemize}
\end{block}

\begin{block}{Universidad de San Buenaventura (Tamayo, 2018)}
\begin{itemize}
    \item Primer CMG nacional para estabilización en reposo
    \item Simulación en ADAMS
    \item Limitación: Sin validación experimental
\end{itemize}
\end{block}

\vspace{0.3cm}
\begin{center}
\fbox{\parbox{0.75\textwidth}{\centering \textit{[Espacio para tabla comparativa]}}}
\end{center}
\end{frame}

\begin{frame}{Antecedentes - Desarrollos Internacionales}
\begin{block}{The Ohio State University (Kalouche, 2014)}
\begin{itemize}
    \item CMG experimental completo
    \item Volante de 4.5 kg a 4000 RPM
    \item Control LQR, estabilización por 30 segundos
\end{itemize}
\end{block}

\begin{block}{American Control Conference (Yetkin et al., 2014)}
\begin{itemize}
    \item Comparación PID vs SMC
    \item PID mejor balance desempeño/eficiencia
    \item SMC mejor respuesta ante perturbaciones
\end{itemize}
\end{block}

\vspace{0.3cm}
\begin{center}
\fbox{\parbox{0.75\textwidth}{\centering \textit{[Espacio para imágenes de prototipos]}}}
\end{center}
\end{frame}

\begin{frame}{Efecto Giroscópico}
\begin{block}{Principio Fundamental}
Cuando un objeto gira rápidamente alrededor de su eje de simetría, posee momento angular $\mathbf{L}$.
\end{block}

\begin{columns}
\column{0.5\textwidth}
\textbf{Momento Angular:}
\begin{equation*}
L = I\omega
\end{equation*}

\textbf{Precesión:}
\begin{equation*}
\Omega_p = \frac{\tau}{L} = \frac{mgr}{I\omega}
\end{equation*}

\column{0.5\textwidth}
\vspace{0.2cm}
\fbox{\parbox{0.9\textwidth}{\centering \textit{[Espacio para diagrama: Efecto giroscópico]}}}
\end{columns}

\vspace{0.3cm}
\begin{alertblock}{Clave del Sistema}
El volante de inercia genera torques de reacción perpendiculares al torque aplicado, compensando la inclinación de la bicicleta.
\end{alertblock}
\end{frame}

\begin{frame}{Controlador PID}
\begin{block}{Ecuación General}
\begin{equation*}
u(t) = K_p e(t) + K_i \int_{0}^{t} e(\tau)d\tau + K_d \frac{de(t)}{dt}
\end{equation*}
\end{block}

\begin{columns}
\column{0.5\textwidth}
\textbf{Componentes:}
\begin{itemize}
    \item $K_p$: Proporcional
    \item $K_i$: Integral
    \item $K_d$: Derivativo
\end{itemize}

\column{0.5\textwidth}
\vspace{0.2cm}
\fbox{\parbox{0.9\textwidth}{\centering \textit{[Espacio para diagrama de bloques PID]}}}
\end{columns}

\vspace{0.3cm}
\begin{block}{Ventajas}
\begin{itemize}
    \item Robustez y simplicidad
    \item Facilidad de sintonización
    \item Bajo costo computacional
\end{itemize}
\end{block}
\end{frame}

\begin{frame}{Sensores y Microcontrolador}
\begin{columns}
\column{0.5\textwidth}
\begin{block}{MPU9250 (IMU 9-DOF)}
\begin{itemize}
    \item Acelerómetro 3 ejes
    \item Giroscopio 3 ejes
    \item Magnetómetro 3 ejes
    \item Sistema MARG
\end{itemize}
\end{block}

\vspace{0.2cm}
\fbox{\parbox{0.9\textwidth}{\centering \textit{[Espacio para imagen: Sensor MPU9250]}}}

\column{0.5\textwidth}
\begin{block}{ESP32}
\begin{itemize}
    \item Dual-core Xtensa 32 bits
    \item Hasta 240 MHz
    \item Procesamiento paralelo
    \item Comunicación I²C
\end{itemize}
\end{block}

\vspace{0.2cm}
\fbox{\parbox{0.9\textwidth}{\centering \textit{[Espacio para imagen: ESP32]}}}
\end{columns}

\vspace{0.3cm}
\begin{alertblock}{Filtro Complementario}
\begin{equation*}
\text{Ángulo} = W \cdot (\text{Ángulo}_{prev} + \omega \cdot \Delta t) + (1-W) \cdot \text{Ángulo}_{accel}
\end{equation*}
\end{alertblock}
\end{frame}

% ============================================
% SECCIÓN 6: MODELO MATEMÁTICO
% ============================================
\section{Modelo Matemático}

\begin{frame}{Modelo Dinámico del Sistema}
\begin{block}{Componentes del Sistema}
\begin{itemize}
    \item \textbf{Bicicleta ($b$):} Estructura principal
    \item \textbf{Gimbal ($g$):} Marco motorizado (eje de precesión)
    \item \textbf{Flywheel ($f$):} Disco de inercia a alta velocidad $\Omega$
\end{itemize}
\end{block}

\vspace{0.2cm}
\begin{center}
\fbox{\parbox{0.8\textwidth}{\centering \textit{[Espacio para diagrama de cuerpo libre]}}}
\end{center}

\begin{block}{Variables de Estado}
\begin{itemize}
    \item $\theta$: Ángulo de inclinación de la bicicleta
    \item $\alpha$: Ángulo de precesión del gimbal
\end{itemize}
\end{block}
\end{frame}

\begin{frame}{Ecuaciones de Movimiento}
\begin{block}{Energía Cinética Total}
\begin{equation*}
T = T_{Bike} + T_{Gimbal} + T_{Flywheel}
\end{equation*}
\end{block}

\begin{block}{Energía Potencial}
\begin{equation*}
V = (m_b h_b + m_g h_g + m_f h_f) g \cos\theta
\end{equation*}
\end{block}

\begin{block}{Ecuación de Inclinación (Linealizada)}
\begin{equation*}
J_{eq}\ddot{\theta} - (M_{eq}L_{eq}g)\theta + H\dot{\alpha} = 0
\end{equation*}
\end{block}

\begin{block}{Ecuación de Precesión}
\begin{equation*}
I_\alpha \ddot{\alpha} - H\dot{\theta} = \tau_\alpha
\end{equation*}
\end{block}
\end{frame}

\begin{frame}{Parámetros del Sistema}
\begin{center}
\fbox{\parbox{0.9\textwidth}{\centering \textit{[Espacio para tabla de parámetros]}}}
\end{center}

\vspace{0.3cm}
\begin{block}{Constantes Importantes}
\begin{itemize}
    \item $J_{eq}$: Inercia efectiva de balanceo
    \item $M_{eq}L_{eq}$: Momento estático total
    \item $I_\alpha$: Inercia total del eje de precesión
    \item $H = I_{fy} \Omega$: Momento angular del flywheel
\end{itemize}
\end{block}
\end{frame}

% ============================================
% SECCIÓN 7: DISEÑO E IMPLEMENTACIÓN
% ============================================
\section{Diseño e Implementación}

\begin{frame}{Criterios de Diseño}
\begin{block}{Decisión de Diseño: Sistema CMG}
\textbf{Razones:}
\begin{enumerate}
    \item Operación a velocidad cero
    \item Validación experimental documentada
    \item Simplicidad conceptual
    \item Control mediante PID
\end{enumerate}
\end{block}

\begin{columns}
\column{0.5\textwidth}
\textbf{Ventajas:}
\begin{itemize}
    \item Estabilización en reposo
    \item Modelado matemático claro
    \item Facilidad de sintonización
\end{itemize}

\column{0.5\textwidth}
\vspace{0.2cm}
\fbox{\parbox{0.9\textwidth}{\centering \textit{[Espacio para esquema CMG]}}}
\end{columns}
\end{frame}

\begin{frame}{Diseño 3D - Prototipo}
\begin{center}
\fbox{\parbox{0.85\textwidth}{\centering \textit{[Espacio para imagen: Diseño CAD completo]}}}
\end{center}

\vspace{0.3cm}
\begin{columns}
\column{0.5\textwidth}
\fbox{\parbox{0.9\textwidth}{\centering \textit{[Espacio para imagen: Vista lateral]}}}

\column{0.5\textwidth}
\fbox{\parbox{0.9\textwidth}{\centering \textit{[Espacio para imagen: Vista frontal]}}}
\end{columns}
\end{frame}

\begin{frame}{Componentes Mecánicos}
\begin{block}{Sistema de Volante de Inercia}
\begin{itemize}
    \item Disco de inercia (flywheel)
    \item Motor DC de alta velocidad
    \item Sistema de transmisión
\end{itemize}
\end{block}

\begin{block}{Sistema Gimbal}
\begin{itemize}
    \item Marco pivotante motorizado
    \item Servomotor de precesión
    \item Rodamientos de baja fricción
\end{itemize}
\end{block}

\vspace{0.3cm}
\begin{center}
\fbox{\parbox{0.8\textwidth}{\centering \textit{[Espacio para imagen: Componentes mecánicos]}}}
\end{center}
\end{frame}

\begin{frame}{Arquitectura del Sistema de Control}
\begin{center}
\fbox{\parbox{0.9\textwidth}{\centering \textit{[Espacio para diagrama: Arquitectura completa]}}}
\end{center}

\vspace{0.3cm}
\begin{block}{Flujo de Control}
\begin{enumerate}
    \item IMU mide inclinación ($\theta$) y velocidad angular
    \item Filtro complementario procesa señales
    \item Controlador PID calcula señal de control
    \item Actuador ajusta ángulo del gimbal ($\alpha$)
    \item Torque giroscópico estabiliza la bicicleta
\end{enumerate}
\end{block}
\end{frame}

\begin{frame}{Sistema Electrónico}
\begin{columns}
\column{0.5\textwidth}
\begin{block}{Componentes}
\begin{itemize}
    \item ESP32 (controlador)
    \item MPU9250 (IMU)
    \item BTS7960 (driver motor)
    \item Batería LiFePO4
    \item BMS (gestión baterías)
\end{itemize}
\end{block}

\column{0.5\textwidth}
\vspace{0.2cm}
\fbox{\parbox{0.9\textwidth}{\centering \textit{[Espacio para diagrama electrónico]}}}
\end{columns}

\vspace{0.3cm}
\begin{center}
\fbox{\parbox{0.8\textwidth}{\centering \textit{[Espacio para imagen: Circuito implementado]}}}
\end{center}
\end{frame}

\begin{frame}{Implementación del Controlador PID}
\begin{block}{Algoritmo de Control}
\begin{itemize}
    \item Frecuencia de muestreo: 100 Hz
    \item Filtro complementario (W = 0.99)
    \item Sintonización de ganancias $K_p$, $K_i$, $K_d$
\end{itemize}
\end{block}

\begin{center}
\fbox{\parbox{0.8\textwidth}{\centering \textit{[Espacio para código: Fragmento algoritmo PID]}}}
\end{center}

\vspace{0.3cm}
\begin{alertblock}{Parámetros Finales}
\begin{itemize}
    \item $K_p$: [valor]
    \item $K_i$: [valor]
    \item $K_d$: [valor]
\end{itemize}
\end{alertblock}
\end{frame}

% ============================================
% SECCIÓN 8: RESULTADOS
% ============================================
\section{Resultados}

\begin{frame}{Prototipo Final}
\begin{center}
\fbox{\parbox{0.85\textwidth}{\centering \textit{[Espacio para imagen: Prototipo completo ensamblado]}}}
\end{center}

\vspace{0.3cm}
\begin{columns}
\column{0.5\textwidth}
\fbox{\parbox{0.9\textwidth}{\centering \textit{[Espacio para imagen: Detalle volante]}}}

\column{0.5\textwidth}
\fbox{\parbox{0.9\textwidth}{\centering \textit{[Espacio para imagen: Detalle gimbal]}}}
\end{columns}
\end{frame}

\begin{frame}{Pruebas de Estabilización}
\begin{block}{Condiciones de Prueba}
\begin{itemize}
    \item Estabilización en reposo
    \item Perturbaciones laterales controladas
    \item Tiempo de estabilización
    \item Ángulo de recuperación
\end{itemize}
\end{block}

\vspace{0.3cm}
\begin{center}
\fbox{\parbox{0.85\textwidth}{\centering \textit{[Espacio para gráfica: Respuesta temporal del sistema]}}}
\end{center}
\end{frame}

\begin{frame}{Resultados Experimentales}
\begin{center}
\fbox{\parbox{0.9\textwidth}{\centering \textit{[Espacio para tabla: Resultados cuantitativos]}}}
\end{center}

\vspace{0.3cm}
\begin{columns}
\column{0.5\textwidth}
\begin{block}{Métricas}
\begin{itemize}
    \item Error en estado estacionario
    \item Tiempo de asentamiento
    \item Sobrepaso máximo
    \item Perturbación máxima
\end{itemize}
\end{block}

\column{0.5\textwidth}
\fbox{\parbox{0.9\textwidth}{\centering \textit{[Espacio para gráfica comparativa]}}}
\end{columns}
\end{frame}

\begin{frame}{Análisis de Desempeño}
\begin{center}
\fbox{\parbox{0.9\textwidth}{\centering \textit{[Espacio para gráfica: Ángulo vs Tiempo]}}}
\end{center}

\vspace{0.3cm}
\begin{columns}
\column{0.5\textwidth}
\fbox{\parbox{0.9\textwidth}{\centering \textit{[Espacio para gráfica: Señal de control]}}}

\column{0.5\textwidth}
\fbox{\parbox{0.9\textwidth}{\centering \textit{[Espacio para gráfica: Error vs Tiempo]}}}
\end{columns}
\end{frame}

\begin{frame}{Validación del Sistema}
\begin{block}{Logros Alcanzados}
\begin{itemize}
    \item \checkmark Estabilización activa en reposo
    \item \checkmark Compensación de perturbaciones
    \item \checkmark Validación del modelo matemático
    \item \checkmark Implementación exitosa del controlador PID
\end{itemize}
\end{block}

\vspace{0.3cm}
\begin{center}
\fbox{\parbox{0.8\textwidth}{\centering \textit{[Espacio para video/secuencia: Demostración funcionamiento]}}}
\end{center}
\end{frame}

% ============================================
% SECCIÓN 9: CONCLUSIONES
% ============================================
\section{Conclusiones}

\begin{frame}{Conclusiones}
\begin{block}{Técnicas}
\begin{itemize}
    \item Se validó experimentalmente el principio de estabilización giroscópica mediante CMG
    \item El controlador PID demostró ser efectivo y eficiente para el objetivo planteado
    \item El modelo matemático captura adecuadamente la dinámica del sistema
\end{itemize}
\end{block}

\begin{block}{Académicas}
\begin{itemize}
    \item Integración exitosa de teoría de control con implementación práctica
    \item Contribución al desarrollo tecnológico nacional en sistemas de control
    \item Generación de conocimiento aplicado a movilidad sostenible
\end{itemize}
\end{block}
\end{frame}

\begin{frame}{Recomendaciones}
\begin{block}{Trabajo Futuro}
\begin{itemize}
    \item Implementar controladores avanzados (ADRC, SMC, LQR)
    \item Optimizar el diseño mecánico para reducir peso
    \item Integrar sistema de control de dirección
    \item Evaluar desempeño en condiciones de movimiento
    \item Desarrollar sistema de navegación autónoma
\end{itemize}
\end{block}

\begin{alertblock}{Aplicaciones Potenciales}
\begin{itemize}
    \item Sistema de asistencia para ciclistas
    \item Plataformas de movilidad personal
    \item Educación en sistemas de control
\end{itemize}
\end{alertblock}
\end{frame}

\begin{frame}{Impacto del Proyecto}
\begin{columns}
\column{0.5\textwidth}
\begin{block}{Social}
\begin{itemize}
    \item Mayor seguridad para ciclistas
    \item Promoción de movilidad sostenible
    \item Inclusión de usuarios con menor destreza
\end{itemize}
\end{block}

\column{0.5\textwidth}
\begin{block}{Ambiental}
\begin{itemize}
    \item Reducción de emisiones
    \item Alternativa de transporte limpio
    \item Contribución a calidad del aire
\end{itemize}
\end{block}
\end{columns}

\vspace{0.5cm}
\begin{block}{Tecnológico}
\begin{itemize}
    \item Avance en control de sistemas dinámicos
    \item Desarrollo de prototipos funcionales
    \item Base para futuras investigaciones
\end{itemize}
\end{block}
\end{frame}

% ============================================
% AGRADECIMIENTOS
% ============================================
\begin{frame}{Agradecimientos}
\begin{center}
\Large
\textbf{Agradecimientos}

\vspace{1cm}
\normalsize

\textbf{Director:} Dr. Alexander Jiménez Triana

\vspace{0.5cm}

\textbf{Grupo de Investigación ORCA}\\
Universidad Distrital Francisco José de Caldas

\vspace{0.5cm}

\textbf{Facultad Tecnológica}\\
Ingeniería en Control y Automatización

\vspace{1cm}

\fbox{\parbox{0.6\textwidth}{\centering \textit{[Espacio para logo Universidad]}}}
\end{center}
\end{frame}

% ============================================
% PREGUNTAS
% ============================================
\begin{frame}
\begin{center}
\Huge
\textbf{¿Preguntas?}

\vspace{2cm}
\normalsize

\textbf{Contacto:}\\
\vspace{0.3cm}
jhhernandezm@udistrital.edu.co\\
fareyesr@udistrital.edu.co

\vspace{1.5cm}

\Large
\textbf{Gracias por su atención}
\end{center}
\end{frame}

% ============================================
% DIAPOSITIVAS DE RESPALDO
% ============================================
\appendix

\begin{frame}{Respaldo: Especificaciones Técnicas}
\begin{center}
\fbox{\parbox{0.9\textwidth}{\centering \textit{[Espacio para tabla: Especificaciones detalladas]}}}
\end{center}

\vspace{0.3cm}
\begin{columns}
\column{0.5\textwidth}
\textbf{Volante de Inercia:}
\begin{itemize}
    \item Masa: [valor]
    \item Momento de inercia: [valor]
    \item Velocidad: [valor] RPM
\end{itemize}

\column{0.5\textwidth}
\textbf{Actuador Gimbal:}
\begin{itemize}
    \item Torque: [valor] N·m
    \item Rango: [valor]°
    \item Velocidad: [valor] rad/s
\end{itemize}
\end{columns}
\end{frame}

\begin{frame}{Respaldo: Diagrama de Bloques Detallado}
\begin{center}
\fbox{\parbox{0.95\textwidth}{\centering \textit{[Espacio para diagrama de bloques completo del sistema de control]}}}
\end{center}
\end{frame}

\begin{frame}{Respaldo: Análisis de Estabilidad}
\begin{center}
\fbox{\parbox{0.9\textwidth}{\centering \textit{[Espacio para lugar de raíces o diagrama de Bode]}}}
\end{center}

\vspace{0.3cm}
\begin{block}{Criterios de Estabilidad}
\begin{itemize}
    \item Márgenes de ganancia y fase
    \item Ubicación de polos en el plano complejo
    \item Análisis de robustez
\end{itemize}
\end{block}
\end{frame}

\begin{frame}{Respaldo: Código del Controlador}
\begin{center}
\fbox{\parbox{0.95\textwidth}{\centering \textit{[Espacio para fragmentos de código del controlador PID y filtro complementario]}}}
\end{center}
\end{frame}

\begin{frame}{Respaldo: Proceso de Construcción}
\begin{columns}
\column{0.5\textwidth}
\fbox{\parbox{0.9\textwidth}{\centering \textit{[Espacio para foto: Fase 1 construcción]}}}

\vspace{0.3cm}
\fbox{\parbox{0.9\textwidth}{\centering \textit{[Espacio para foto: Fase 2 construcción]}}}

\column{0.5\textwidth}
\fbox{\parbox{0.9\textwidth}{\centering \textit{[Espacio para foto: Fase 3 construcción]}}}

\vspace{0.3cm}
\fbox{\parbox{0.9\textwidth}{\centering \textit{[Espacio para foto: Prototipo final]}}}
\end{columns}
\end{frame}

\begin{frame}{Respaldo: Comparación con Estado del Arte}
\begin{center}
\fbox{\parbox{0.95\textwidth}{\centering \textit{[Espacio para tabla comparativa con otros proyectos]}}}
\end{center}

\vspace{0.3cm}
\begin{block}{Contribuciones Propias}
\begin{itemize}
    \item Primera implementación física nacional de CMG
    \item Validación experimental completa
    \item Plataforma de bajo costo
\end{itemize}
\end{block}
\end{frame}

\end{document}
